\begin{resumo}[Abstract]
  \begin{otherlanguage*}{english}

This work describes the activities carried out by student Júlia Nahas
Ribeiro de Araújo during the supervised internship in the Equine Medical
Clinic sector of the Veterinary Hospital of the Federal University of
Minas Gerais (UFMG), located in Belo Horizonte, Minas Gerais, Brazil, in
fulfillment of the requirements of discipline PRG -- 107 Supervised
Internship. This discipline is part of the tenth-period curriculum of
the Veterinary Medicine undergraduate program and comprises 476 hours,
divided between practical and theoretical activities. The internship was
conducted from March 9 to May 22, 2026, under the supervision of Profa.
Dra Ana Luísa Soares de Miranda and guidance of Profa. Dra Ana Paula
Peconick, totaling 424 practical hours. Throughout the internship, the
student participated in the clinical routine of the sector and in
external consultations, attending appointments, diagnostic examinations
and clinical procedures across different areas of equine medicine, as
well as clinical case discussions and other complementary activities.
These experiences contributed to consolidating the integration between
the theoretical knowledge acquired throughout the undergraduate program
and professional practice. This work describes, in sequence, the
infrastructure and operation of the sector, the caseload recorded during
the period and, finally, a clinical case report followed during the
supervised internship.

\vspace{\onelineskip}

\noindent
\textbf{Keywords}: Supervised internship. Equine medicine. Rhodococcosis.

  \end{otherlanguage*}
\end{resumo}
